% ECONOMICS_PROBLEM: {"id": "O32-614", "title": "Organizational complements", "jel_code": "O32", "source_id": "OP-0414"}
% !TeX program = xelatex
\documentclass[11pt,a4paper]{article}
\usepackage[margin=23mm]{geometry}
\usepackage{fontspec}
\setmainfont{Latin Modern Roman}
\usepackage{amsmath,amssymb,mathtools}
\usepackage{xcolor,enumitem,booktabs,longtable,array,xurl,hyperref}
\definecolor{muted}{HTML}{53636C}
\hypersetup{colorlinks=true,urlcolor=blue,linkcolor=blue}
\setlength{\parindent}{0pt}
\setlength{\parskip}{5pt}
\newcommand{\R}{\mathbb R}
\newcommand{\E}{\mathbb E}
\newcommand{\N}{\mathbb N}
\newcommand{\ind}{\mathbf 1}
\newcommand{\Var}{\operatorname{Var}}
\newcommand{\Cov}{\operatorname{Cov}}
\newcommand{\supp}{\operatorname{supp}}
\DeclareMathOperator*{\argmin}{arg\,min}
\DeclareMathOperator*{\argmax}{arg\,max}
\newcommand{\entrymeta}[3]{{\sffamily\small\color{muted}\textbf{\texttt{#1}}\quad|\quad #2\par #3\par}}
\newcommand{\Problemref}[1]{\texttt{#1}}
\newcommand{\JELref}[2]{\texttt{#2} (JEL \texttt{#1})}
\begin{document}
\section*{Organizational complements}
\textbf{Problem ID:} \texttt{O32-614}\quad \textbf{JEL:} \texttt{O32}\par
% BEGIN ECONOMICS STATEMENT
\label{problem:OP-0414}
{\sffamily\small\color{muted}Original subject group: Innovation and AI economics\par}
\label{definitions:problem:30007696}
\textbf{Audit classification:} Background; proposed extension. The cited work supports AI and R\&D reorganization as background. No exact published open statement of this menu-constrained team design was verified.
\subsection*{Definitions and precise research target}
Consider research teams inside firms conducting comparable product-development projects, with a fixed two-year validation horizon. Let \(A\in\{0,1\}\) indicate access to an AI research assistant, and let \(Z\in\mathcal Z\) denote one of finitely many preregistered organizational designs. Each design specifies task allocation, experiment approval rules, independent verification responsibilities and the distribution of decision authority. Design \(z_0\) is the existing organization. Let \(Y_i(a,z)\) be independently validated discovery value and \(C_i(a,z)>0\) the real labor, computation, experimentation and review cost for team \(i\). Define \(P(a,z)=E[Y_i(a,z)]/E[C_i(a,z)]\) over eligible research teams, and the organizational complementarity
\[
\Delta(z)=P(1,z)-P(0,z)-P(1,z_0)+P(0,z_0).
\]
Randomize AI access and organizational designs across sufficiently isolated teams, or across connected team clusters where knowledge sharing occurs. Hold project eligibility and evaluation criteria fixed, include failed projects, and distinguish faster drafting from reproducible technical discoveries. The design problem is to choose \(z^\ast\in\arg\max_{z\in\mathcal Z}P(1,z)\), subject to stated safety and review-capacity constraints. Which redistribution of research work and authority creates positive complementarity with AI, and what mechanisms explain it? This specification concerns AI-assisted R\&D organization rather than economy-wide adoption delays or aggregate intangible investment.

\textbf{Definitions, assumptions and mathematical qualifications.} The organizational menu \(\mathcal Z\) is finite and each element fixes project allocation, authority, approval and verification rules. Require \(0<c_{\min}\leq E[C_i(a,z)]<\infty\) and finite nonnegative discovery value, measured consistently in both arms. \(P(a,z)\) is a ratio of means, not mean individual productivity; its interaction need not equal an interaction in gross value. A feasible design satisfies explicit conditions such as expected unsafe failures at most \(\rho\) and expected verification hours at most \(K\). Require a nonempty feasible subset before defining \(z^*\). Include within-team interactions in the treatment rather than assume that individual task effects add.
\subsection*{Variants and assumption changes}
\begin{enumerate}
\item \textbf{Editorial, productivity.} Optimize the given ratio \(P(1,z)\) under explicit verification-capacity constraints; report both its numerator and denominator.
\item \textbf{Editorial, monetary net value.} Optimize \(E[Y_i(1,z)-C_i(1,z)]\), permitting project reassignment but holding the eligible project cohort fixed; productivity and net-value maximizers may differ.
\end{enumerate}
\subsection*{References and source audit}
\textbf{Checked source.} NBER Working Paper 24449 (2018), indexed primary abstract. The publication motivates AI as a method of invention and R\&D reorganization. Primary metadata and abstract indexed; exact organization-design agenda passage not retrieved. \url{https://www.nber.org/papers/w24449}.
\begin{enumerate}\raggedright
\item\hypertarget{definitions:source:30007696:1}{} Iain M. Cockburn; Rebecca Henderson; Scott Stern. 2018. \emph{The Impact of Artificial Intelligence on Innovation}. NBER Working Paper 24449 (2018), indexed primary abstract.
\url{https://www.nber.org/papers/w24449}.
DOI: \url{https://doi.org/10.3386/w24449}.
\end{enumerate}

\subsection*{Common conventions: Source mathematical conventions}

These conventions apply to each entry unless explicitly replaced.

\textbf{Models and identification.} A primitive model \(m\) specifies all probability laws, feasible choices, information, equilibrium and measurement rules used by its target. The admissible class \(\mathcal M\) is fixed independently of outcomes. If \(P_O(m)\) is its observable probability law and \(\tau(m)\) the target, its identified set at observable law \(P\) is
\[
 \mathcal I_\tau(P)=\{\tau(m):m\in\mathcal M,\ P_O(m)=P\}.
\]
Point identification means that this set is a singleton. An empty set signals incompatibility of data and assumptions. Coordinate bounds need not describe the joint set. Bounds are sharp only if every retained target is attainable or arbitrarily approachable by compatible models. Finite bounds require explicit restrictions on unobserved quantities; unrestricted confounding, hidden wealth, extrapolation or arbitrary equilibrium selection can yield uninformative sets.

\textbf{Causal interventions.} A potential outcome \(Y_i(p)\) is the outcome under a complete policy regime \(p\), including timing, financing, information and market spillovers. The target population and units are fixed before treatment. With interference, use the complete assignment vector or a justified exposure mapping; individual no-interference notation alone cannot identify an economy-wide intervention. A difference of mean potential outcomes does not require the cross-world joint distribution. Conditioning on post-treatment migration, survival, enrollment or employment changes the estimand and needs separate justification.

\textbf{Statistical statements.} An estimator, confidence set, forecast or test specifies a sampling law, model class and loss/coverage criterion. Uniform \((1-\alpha)\)-coverage means \(\inf_{m\in\mathcal M}\Pr_m\{\tau(m)\in C_n\}\geq1-\alpha\), or an explicitly asymptotic version. The whole parameter space trivially has coverage, so informativeness requires a stated width, risk or power objective. A forecasting improvement is relative to a named benchmark, information set and evaluation population.

\textbf{Equilibrium and optimization.} An equilibrium comprises feasible plans, optimality, beliefs where needed, budget constraints and all market/resource clearing. The primitives or a stated benchmark define each correspondence \(\mathcal E(p;m)\). Nonemptiness is assumed or proved, never inferred just from compact policy sets. A selected-equilibrium target uses a declared selection rule; a robust target quantifies over all permitted equilibria. Use suprema or infima unless attainment is established. An \(\varepsilon\)-optimal policy is feasible and within \(\varepsilon\) of the finite optimum value. Report an empty feasible set separately.

\textbf{Welfare and accounting.} Normative utility, interpersonal normalization, social weights, numeraire, discounting and the use of public receipts are primitives. Behavioral utility need not coincide with the welfare criterion. Household welfare includes ownership income and rents once; adding profits again to compensating variation generally double-counts them. A monetary equivalent is defined by an explicit transfer and price convention, with monotonicity and range conditions. Average effects, mechanism contributions, transfers and welfare losses are different targets. Infinite sums require convergence and dynamic feasible plans require terminal or no-Ponzi conditions.

\textbf{Source versions.} A working-paper date, revision date and journal publication year are distinguished. A DOI for a journal version is not automatically the DOI of an earlier manuscript. Locators refer to the stated version and printed pages where specified. A metadata-only check does not verify a claim about a full-text conclusion. Every entry's source subsection narrows the provenance claim accordingly.
% END ECONOMICS STATEMENT
\end{document}
